%*******************************************************************
% Chapter 2: Related Work
% UniAdapFuse: Reliability-Aware Multimodal Fusion for
% UAV Disaster Perception
%*******************************************************************

\section{Related Work}

This chapter surveys prior research across the foundational pillars of our framework: (1) single-modality aerial sensing and physical failure modes; (2) multimodal fusion architectures and 2024--2026 contemporary SOTA paradigms; (3) real-time aerial object detection and vision transformers; (4) multi-task negative transfer in unified perception; and (5) edge-deployment compression on UAV embedded hardware. For each domain, we analyze existing methodologies, identify critical failure modes, and establish the technical rationale for the architectural innovations in \textbf{UniAdapFuse}.

\subsection{Single-Modality Aerial Sensing and Environmental Failure Modes}

The three primary modalities used in aerial UAV disaster response (RGB visual cameras, thermal infrared sensors, and acoustic microphones) capture fundamentally distinct physical phenomena, exhibiting complementary strengths and misaligned failure modes.

\textbf{RGB Visual Sensing.}
Visible-light cameras provide high spatial resolution (typically $1920\times1080$ to $4\text{K}$) and rich chromatic signatures essential for distinguishing flame color palettes, smoke plume textures, and human survivor clothing~\cite{celik2009fire, toreyin2006computer}. Deep convolutional neural networks (CNNs) have achieved high image-level fire detection and localization accuracy in video surveillance settings~\cite{muhammad2019efficient}. However, RGB systems suffer from catastrophic physical vulnerability: dense wildfire smoke reduces visibility by 50--90\%, nighttime flight eliminates chromatic information, and solar glare saturates camera sensors. In dense smoke, RGB fire detection accuracy collapses to 20--40\%, while search-and-rescue victim localization drops below 15\%. These failure modes cannot be resolved by scaling RGB model parameters alone; they represent fundamental physical limits of optical sensing.

\textbf{Thermal Infrared Imaging.}
Thermal infrared sensors detect long-wave infrared radiation (LWIR, 7.5--14 $\mu$m) emitted as blackbody radiation proportional to object temperature. Thermal imaging operates independently of ambient illumination, penetrating darkness, light smoke, and foliage haze to detect active combustion zones and $37^\circ\text{C}$ living human body heat~\cite{toulouse2017computer}. However, aerial thermal sensing carries severe operational limitations: typical UAV thermal sensors have low spatial resolution ($320\times240$ px), preventing distant survivor detection from high altitudes ($>50$ m). Furthermore, solar heating of asphalt, dry soil, and corrugated metal roofs frequently reaches 60--80$^\circ\text{C}$ in daytime flight, triggering severe false-positive detections. In fusion models with fixed modality weights, these solar thermal artifacts can mislead the network and produce excessive false alarms.

\textbf{Acoustic Sensing and UAV Rotor Noise.}
Acoustic sensors capture omnidirectional pressure waves that propagate around visual occlusions, providing distinctive spectral signatures for roaring flames (500--5000 Hz), structural collapse transients ($<500$ Hz), and human distress calls. However, UAV rotor blades spinning at 4000--8000 RPM generate intense acoustic contamination (80--100 dB SPL) that severely degrades signal-to-noise ratios (SNR). Standard acoustic classifiers without dynamic reliability estimation fail under rotor noise, necessitating autonomous gating mechanisms that suppress acoustic tokens when noise dominates.

\subsection{Real-Time Aerial Object Detection and Vision Transformers}

Autonomous UAV disaster response requires \textbf{dense 2D bounding box localization} $(\text{xmin}, \text{ymin}, \text{xmax}, \text{ymax})$ to direct tactical fire suppression and rescue operations.

\textbf{Anchor-Free Convolutional Detectors (YOLO Series).}
The You Only Look Once (YOLO) family represents the standard for real-time edge detection. Recent advancements including YOLOv8, YOLOv10~\cite{wang2024yolov10}, YOLO11~\cite{ultralytics2024yolo11}, and YOLO26~\cite{sapkota2025yolo26} have adopted anchor-free heads and C3k2 cross-stage feature blocks, and YOLOv10 and YOLO26 remove the Non-Maximum Suppression (NMS) step through dual-label assignment. On the D-Fire smoke and fire benchmark~\cite{dfire2022}, the lightweight real-time detectors trained in this work reach 67.7\% to 72.2\% $\text{mAP}_{50}$. However, standard YOLO architectures rely entirely on data-driven convolutional filters and lack physical domain priors, which may slow learning on diffuse, boundary-less smoke plumes (the baseline YOLO26n run in this work trains for 4.45 hours). UniAdapFuse includes a \textbf{Hazard Prior Gate (HPG)}, which injects explicit chromaticity ($2R - G - B$) and edge residuals into the detection backbone at stride 4; in our runs it reduces wall-clock training time by $2.27\times$ (1.96 versus 4.45 hours) at a cost of 0.94 points of $\mathrm{mAP}_{50}$, although a convergence-rate improvement is not established without learning curves.

\textbf{Real-Time Detection Transformers (RT-DETR).}
To eliminate handcrafted NMS post-processing, RT-DETR~\cite{zhao2024rtdetr} introduces deformable multi-scale cross-attention with bipartite Hungarian matching. While the RT-DETRv2-R18 variant reaches high detection accuracy ($73.8\%\, \text{mAP}_{50}$ on D-Fire in this work), its multi-scale self-attention modules demand high GPU memory and power draw ($>20$W), exceeding the payload constraints of micro-UAV platforms.

\subsection{Public Benchmarks for Aerial Disaster Classification and Fire Detection}
\label{sec:public_benchmarks_related}

Three public datasets used in this thesis have published protocols against which single-modality results can be compared.

\textbf{AIDER.} Kyrkou and Theocharides introduced the AIDER aerial emergency dataset with a lightweight classifier~\cite{kyrkou2019aider} and later EmergencyNet, which uses atrous convolutional feature fusion to classify drone imagery at low cost~\cite{kyrkou2020emergencynet}. Lee et al.\ re-trained six lightweight CNNs from scratch on the 6{,}433-image AIDER subset with the 4:1:2 train/validation/test protocol (2{,}473 test images, class-imbalanced towards the normal class) and proposed WATT-EffNet, a wider, attention-augmented EfficientNet under 1\,M parameters~\cite{lee2023wattefnet}. Under that protocol they report F1 scores of 80.0\% for EfficientNet, 81.0\% for MobileNetV2, 83.1\% for EmergencyNet and 88.5\% for WATT-EffNet.

\textbf{FLAME.} Shamsoshoara et al.\ released FLAME with a fixed fire/no-fire classification split in which the 39{,}375 training frames come from one drone and camera (Matrice 200, Zenmuse X4S) and the 8{,}617 test frames from another (Phantom 3). Their Xception classifier reaches 76.23\% test accuracy~\cite{shamsoshoara2021aerial}. Because the camera changes between training and test, this protocol measures cross-camera generalisation, which a random frame-level split does not.

\textbf{D-Fire.} De Ven\^{a}ncio et al.\ built D-Fire (21{,}527 images with smoke and fire boxes) and showed that filter pruning can cut YOLOv4's computational cost by up to 83.60\% for low-power fire detection~\cite{dfire2022}; the pruned detector is reported to keep the baseline's 73.98\% mAP.

These works are single-modality and evaluate each dataset under its own protocol. Section~\ref{sec:public_benchmark_results} trains our models under the AIDER and FLAME protocols so that they can be compared with these published results.

\subsection{Contemporary 2024--2026 Multimodal SOTA Literature}

Multimodal fusion for aerial robotics has advanced rapidly from static concatenation to dynamic gating:

\begin{itemize}
    \item \textbf{UAV RGB-Thermal Fusion Dataset and Model (Zhang et al., 2025)~\cite{zhang2025rgbt}:} A UAV-collected multi-scenario RGB-thermal dataset paired with a fusion model for enhanced forest fire detection. While effective for paired optical-thermal imagery, it lacks acoustic integration, physical hazard priors, and nuisance false-alarm modeling.
    \item \textbf{Multi-Modal Training Dynamics (Wang et al., 2020)~\cite{wang2020makes}:} Wang et al.~\cite{wang2020makes} investigated the challenges of training multi-modal classification networks, including modality-specific overfitting and generalization. However, their analysis does not address UAV-specific rotor noise, thermal sensor drift, or explicit physical nuisance modeling.
    \item \textbf{DeepFire (Khan et al., 2022)~\cite{khan2022deepfire}:} DeepFire~\cite{khan2022deepfire} provides a dataset and deep transfer-learning benchmark for forest-fire detection. It is RGB-only and therefore offers no mechanism to down-weight a blinded camera during heavy smoke.
    \item \textbf{DroneAudioset (Gupta et al., 2025)~\cite{gupta2025droneaudioset}:} A 23.5-hour synchronized UAV SAR acoustic dataset with measured rotor-noise profiles, providing the empirical foundation for acoustic search-and-rescue and noise-resilient multimodal fusion.
\end{itemize}

\subsection{Multi-Task Learning and Negative Transfer in Unified Vision}

Simultaneous prediction of global scene semantics (disaster categorization, sensor trust) and dense spatial objects (bounding box regression) within a single shared backbone is notoriously vulnerable to \textbf{multi-task negative transfer}~\cite{yu2020gradient}. When dense regression losses (CIoU and objectness) are backpropagated concurrently with classification losses, high-magnitude localization gradients disrupt the shared feature representations, causing classification accuracy to collapse (as observed in our naive UniAdapFuse v1 baseline, where disaster accuracy dropped from 99.4\% to 53.7\%).

Prior multi-task works attempt to mitigate this through gradient surgery (e.g., PCGrad) or loss weighting. In this work, phased curriculum training with a 15-epoch classification warmup and stop-gradient isolation (\texttt{.detach()}) on multi-scale feature pyramids recovers scene classification to 98.57\% Combined F1. The recorded detection $\mathrm{mAP}_{50}$ for that joint checkpoint is 0.0, so the result demonstrates classification recovery but not successful joint localization.

\subsection{Edge Deployment and Knowledge Distillation}

Embedded UAV compute platforms (typically constrained to strict 10--15W power envelopes) impose severe memory and latency constraints. While structured pruning~\cite{zhu2017prune} and INT8 quantization~\cite{jacob2018quantization} compress individual backbones, compressing multi-branch multimodal networks requires preserving cross-modal alignment. Knowledge distillation (KD)~\cite{hinton2015distilling} transfers rich relational features from a high-capacity teacher to a compact student. In our framework, AdapFuse-v2 leverages feature and logit KD to halve the shared projection width (192 $\to$ 96). Because the two dominant visual backbones (ResNet-18, MobileNetV3-Small) are untouched by this change, the measured total parameter reduction is modest, 12.94M $\to$ 12.30M (4.9\%), rather than a proportional halving; the teacher AdapFuse-v1 measures 70.1 FPS (14.26 ms) on the evaluation workstation GPU, and dedicated edge-hardware profiling of the distilled student is left to future work (see Chapter~6).

\subsection{Research Gaps and Synthesis}

Table~\ref{tab:sota_gap_map} synthesizes the identified literature gaps and maps them to our architectural solutions.

\begin{table}[H]
\centering
\caption{Synthesis of Literature Gaps and UniAdapFuse Architectural Responses}
\label{tab:sota_gap_map}
\resizebox{\textwidth}{!}{%
\begin{tabular}{p{4.2cm} p{5.5cm} p{5.5cm}}
\toprule
\textbf{Literature Gap} & \textbf{UniAdapFuse Architectural Response} & \textbf{Observed Empirical Outcome} \\
\midrule
High UAV detector training cost & \textbf{Hazard Prior Gate (HPG):} Injects physical chromaticity ($2R-G-B$) and edge residuals at stride 4 & $70.53\%\, \text{mAP}_{50}$ in \textbf{1.96h} vs.\ $71.47\%$ in 4.45h for YOLO26n ($2.27\times$ shorter run, $-0.94$ pp); convergence rate not tested \\
\addlinespace
Multi-task negative transfer in joint perception & \textbf{Curriculum Warmup \& Stop-Gradient Isolation:} Detaches pyramid features before detection neck & Disaster accuracy recovered from $53.71\% \to 99.26\%$; joint detection $\mathrm{mAP}_{50}=0.0$ \\
\addlinespace
Static fusion weights under camera failure & \textbf{Reliability \& Uncertainty Estimator (RUE):} Learned scalar gates $r_i \in [0,1]$ & Smoke-overlay F1 is 89.7\% at severity 5, but saved gates remain nearly constant (no dynamic reweighting shown), no matched static baseline is logged, and fused F1 (98.74\%) does not exceed RGB-only (98.85\%) \\
\addlinespace
High operational false alarm rates (12--18\%) & \textbf{5-Class NuisanceAwareHead:} Defines logits for clean signal, solar heating, false fire, audio alarm, occlusion & Branch is unsupervised in our run; 0.73\% is a row-level classifier FPR, not a nuisance-head effect \\
\addlinespace
Edge compute bottleneck on embedded UAVs & \textbf{Knowledge Distillation (AdapFuse-v2):} Halves projection width via logit \& feature KD & 12.30M params (4.9\% smaller) with $<0.2$ pp F1 loss; teacher runs at 70.1 FPS (14.26 ms), edge-GPU profiling of the student left to future work \\
\bottomrule
\end{tabular}%
}
\end{table}
